% META: {"id":"1SPE-PROBCOND-EV-A","chapitre":"1SPE-PROBA-COND","type_objet":"evaluation","version":"A","duree_min":55,"bareme_total":20,"capacites_codes":["C1","C2","C3","C4","C5"],"competences":["chercher","modeliser","representer","calculer","raisonner","communiquer"],"mode_creation":"ex_nihilo","sources_inspiration":[],"parametres_sympy":{},"fichier_tex":"chapitres/1SPE-PROBA-COND/evaluations/1SPE-PROBCOND-EV-A.tex","corrige_tex":"chapitres/1SPE-PROBA-COND/evaluations/1SPE-PROBCOND-EV-A-corrige.tex","status":"approved","origin":"generated","status_history":["generated","audited","accepted","approved"],"release_acceptance":"RELEASE_OWNER_BATCH_ACCEPTANCE","acceptance_closure_digest":"sha256:9b3ccf9a81c5520fbb7e03b7b2d3e7bf2057a02834b6d908bf3be5f49f3b553f","acceptance_receipt_digest":"sha256:4fad86f0282e0fa4e0419cca1ad6379ae7af5a280c04a4a90880f90a1c044242"}
% BEGIN-VERIFY
% from sympy import Rational, simplify
%
% # ---- Exercice 1 : proba conditionnelle + arbre (C1, C2) ----
% P_R = Rational(3, 5)
% P_V = Rational(2, 5)
% P_G_R = Rational(1, 2)
% P_G_V = Rational(1, 4)
% P_RinterG = P_R * P_G_R
% assert P_RinterG == Rational(3, 10)
% P_VinterG = P_V * P_G_V
% assert P_VinterG == Rational(1, 10)
% P_G = P_RinterG + P_VinterG
% assert P_G == Rational(2, 5)
%
% # ---- Exercice 2 : prob totales + Bayes (C3, C5) ----
% P_M = Rational(3, 100)
% P_T_M = Rational(9, 10)
% P_T_Mbar = Rational(1, 25)
% P_T = P_M * P_T_M + (1-P_M) * P_T_Mbar
% assert P_T == Rational(3,100)*Rational(9,10) + Rational(97,100)*Rational(1,25)
% assert P_T == Rational(27,1000) + Rational(97,2500)
% assert P_T == Rational(27,1000) + Rational(388,10000)
% # 27/1000 = 270/10000, so P_T = (270+388)/10000 = 658/10000 = 329/5000
% assert P_T == Rational(270+388, 10000)
% assert P_T == Rational(329, 5000)
% P_M_T = P_M * P_T_M / P_T
% assert P_M_T == Rational(27,1000) / Rational(329,5000)
% assert P_M_T == Rational(27,1000) * Rational(5000,329)
% assert P_M_T == Rational(135, 329)
%
% # ---- Exercice 3 : independance (C4) ----
% P_A = Rational(1, 2)
% P_B = Rational(1, 5)
% P_AinterB = Rational(1, 10)
% assert P_A * P_B == P_AinterB  # independants
% P_AunionB = P_A + P_B - P_AinterB
% assert P_AunionB == Rational(3, 5)
%
% # ---- Exercice 4 : probleme contextualise assurance (C3, C5) ----
% P_J = Rational(1, 5)
% P_S = Rational(4, 5)
% P_A_J = Rational(1, 10)
% P_A_S = Rational(1, 40)
% P_A = P_J * P_A_J + P_S * P_A_S
% assert P_A == Rational(1,5)*Rational(1,10) + Rational(4,5)*Rational(1,40)
% assert P_A == Rational(1,50) + Rational(1,50)
% assert P_A == Rational(1, 25)
% P_J_A = P_J * P_A_J / P_A
% assert P_J_A == Rational(1,50) / Rational(1,25)
% assert P_J_A == Rational(1, 2)
% END-VERIFY

%---------------------------------------------------------------
% Duree : 55 minutes --- Bareme : 20 points
% Toute reponse doit être justifiee.
%---------------------------------------------------------------

\begin{evaluation}{1SPE-PROBCOND-EV-A}{55}

%===============================================================
\section*{Exercice 1 \hfill (5 points) \hfill Capacites C1, C2}
%===============================================================

\textit{Competences evaluees : calculer, représenter.}

\medskip
Un sac contient $3$ jetons rouges et $2$ jetons verts. On tire un jeton au hasard. Si le jeton est rouge, on gagne avec probabilité $\dfrac{1}{2}$. Si le jeton est vert, on gagne avec probabilité $\dfrac{1}{4}$.

\begin{enumerate}
  \item Construire l'arbre pondere de cette expérience.
  \hfill\textit{(1,5 pt --- représenter)}

  \item Calculer $P(R \cap G)$ ou $G$ designe << gagner >>.
  \hfill\textit{(1 pt --- calculer)}

  \item Calculer la probabilité de gagner.
  \hfill\textit{(1,5 pt --- calculer)}

  \item Sachant qu'on a gagne, quelle est la probabilité d'avoir tire un jeton rouge ?
  \hfill\textit{(1 pt --- raisonner)}
\end{enumerate}

%===============================================================
\section*{Exercice 2 \hfill (6 points) \hfill Capacites C3, C5}
%===============================================================

\textit{Competences evaluees : modeliser, calculer, communiquer.}

\medskip
Un test de depistage a une sensibilite de $90\,\%$ et une specificite de $96\,\%$. La prevalence de la maladie est de $3\,\%$. On note $M$ : << malade >> et $T$ : << test positif >>.

\begin{enumerate}
  \item Donner les valeurs de $P(M)$, $P_M(T)$ et $P_{\overline{M}}(T)$.
  \hfill\textit{(1,5 pt --- modeliser)}

  \item Calculer $P(T)$ par la formule des probabilités totales.
  \hfill\textit{(2 pts --- calculer)}

  \item Calculer la valeur predictive positive $P_T(M)$.
  \hfill\textit{(1,5 pt --- calculer)}

  \item Le test est-il fiable pour un diagnostic individuel ? Argumenter.
  \hfill\textit{(1 pt --- communiquer)}
\end{enumerate}

%===============================================================
\section*{Exercice 3 \hfill (4 points) \hfill Capacité C4}
%===============================================================

\textit{Competences evaluees : raisonner, calculer.}

\medskip
On donne $P(A) = \dfrac{1}{2}$, $P(B) = \dfrac{1}{5}$ et $P(A \cap B) = \dfrac{1}{10}$.

\begin{enumerate}
  \item Les événements $A$ et $B$ sont-ils indépendants ? Justifier.
  \hfill\textit{(2 pts --- raisonner)}

  \item Calculer $P(A \cup B)$.
  \hfill\textit{(1 pt --- calculer)}

  \item Montrer que $\overline{A}$ et $B$ sont aussi indépendants.
  \hfill\textit{(1 pt --- raisonner)}
\end{enumerate}

%===============================================================
\section*{Exercice 4 \hfill (5 points) \hfill Capacites C3, C5}
%===============================================================

\textit{Competences evaluees : modeliser, calculer, raisonner.}

\medskip
Une compagnie d'assurance distingue les conducteurs << jeunes >> ($20\,\%$) et << seniors >> ($80\,\%$). La probabilité d'accident est $\dfrac{1}{10}$ pour un jeune et $\dfrac{1}{40}$ pour un senior.

\begin{enumerate}
  \item Construire un arbre pondere.
  \hfill\textit{(1 pt --- représenter)}

  \item Calculer la probabilité qu'un conducteur ait un accident.
  \hfill\textit{(2 pts --- calculer)}

  \item Un conducteur a eu un accident. Calculer la probabilité qu'il soit jeune.
  \hfill\textit{(2 pts --- raisonner)}
\end{enumerate}

\end{evaluation}
